% =====================================================================
%  preamble.tex -- packages, page layout, and box environments
%  (Notation macros live in macros.tex.)
% =====================================================================

% ---------------------------- Fonts & encoding ------------------------
\usepackage[utf8]{inputenc}
\usepackage[T1]{fontenc}
\usepackage{lmodern}
\usepackage{microtype}

% ---------------------------- Page layout -----------------------------
\usepackage[letterpaper, top=1in, bottom=1in, left=1in, right=1.35in,
            marginparwidth=0.95in, marginparsep=0.15in]{geometry}
\setlength{\parindent}{0pt}
\setlength{\parskip}{0.4em}

% ---------------------------- Math ------------------------------------
\usepackage{amsmath, amssymb, amsthm, mathtools}
\usepackage{bm}

% ---------------------------- General ---------------------------------
\usepackage{etoolbox}
\usepackage{xparse}
\usepackage[dvipsnames, svgnames]{xcolor}
\usepackage{graphicx}
\usepackage{booktabs, tabularx, array, multirow}
\usepackage[shortlabels]{enumitem}
\setlist{itemsep=0.1em, topsep=0.25em, parsep=0pt}
\setlist[itemize,2]{label=\textendash}
\setlist[itemize,3]{label=\textperiodcentered}
\usepackage{caption, subcaption}
\captionsetup{font=small, labelfont={bf,sf}}
\usepackage{marginnote}
\usepackage{algorithm}
\usepackage[noend]{algpseudocode}

% ---------------------------- Figures ---------------------------------
\usepackage{tikz}
\usetikzlibrary{arrows.meta, positioning, shapes.geometric, shapes.misc,
                calc, fit, backgrounds, decorations.pathreplacing,
                matrix, quotes, patterns}
\usepackage{pgfplots}
\pgfplotsset{compat=1.18}

% Shared TikZ styles for graph drawings (use these for consistency)
\tikzset{
  vertex/.style   = {circle, draw=black!70, fill=white, thick,
                     minimum size=6.5mm, inner sep=0pt, font=\small},
  vtx/.style      = {circle, draw=black!70, fill=black!10, thick,
                     minimum size=3.5mm, inner sep=0pt},
  target/.style   = {vertex, fill=YellowOrange!45},
  hl/.style       = {vertex, fill=SkyBlue!45},
  edge/.style     = {thick, black!60},
  dedge/.style    = {thick, black!60, -{Stealth[length=2.2mm]}},
  arrow/.style    = {thick, -{Stealth[length=2.6mm]}},
  block/.style    = {draw, rounded corners=2pt, thick, fill=black!4,
                     align=center, inner sep=4pt, font=\small},
  nnbox/.style    = {draw, rounded corners=1pt, thick, fill=Gray!30,
                     minimum width=5mm, minimum height=5mm},
}

% ---------------------------- Colours ---------------------------------
\definecolor{keyblue}{HTML}{1F5FA8}
\definecolor{intuitionpurple}{HTML}{6A3D9A}
\definecolor{cautionred}{HTML}{B2182B}
\definecolor{examorange}{HTML}{D35400}
\definecolor{connectgreen}{HTML}{1B7837}
\definecolor{examplegray}{HTML}{4D4D4D}
\definecolor{asideteal}{HTML}{00796B}
\definecolor{answercolor}{HTML}{1B5E20}

% ---------------------------- Hyperlinks ------------------------------
\usepackage[colorlinks=true, linkcolor=keyblue!80!black,
            urlcolor=keyblue, citecolor=connectgreen,
            bookmarksnumbered=true]{hyperref}
\usepackage[capitalise, noabbrev]{cleveref}
\crefname{chapter}{Lecture}{Lectures}
\Crefname{chapter}{Lecture}{Lectures}

% ---------------------------- Headers ---------------------------------
\usepackage{fancyhdr}
\pagestyle{fancy}
\fancyhf{}
\fancyhead[L]{\small\sffamily\nouppercase{\leftmark}}
\fancyhead[R]{\small\sffamily CS224W Notes}
\fancyfoot[C]{\small\thepage}
\renewcommand{\headrulewidth}{0.4pt}
\setlength{\headheight}{14pt}
\renewcommand{\chaptermark}[1]{\markboth{Lecture \thechapter: #1}{}}

% Chapters are lectures
\renewcommand{\chaptername}{Lecture}
\usepackage{titlesec}
\titleformat{\chapter}[display]
  {\normalfont\sffamily\bfseries\raggedright}
  {\large\color{keyblue}\MakeUppercase{\chaptername}~\thechapter}
  {4pt}{\huge}
\titlespacing*{\chapter}{0pt}{-24pt}{10pt}
\titleformat{\section}{\normalfont\Large\bfseries}{\thesection}{0.8em}{}
\titleformat{\subsection}{\normalfont\large\bfseries}{\thesubsection}{0.8em}{}
\titlespacing*{\section}{0pt}{1.3em}{0.5em}
\titlespacing*{\subsection}{0pt}{1.0em}{0.35em}
\titlespacing*{\paragraph}{0pt}{0.6em}{0.8em}

% ---------------------------- Theorem-like ----------------------------
\usepackage[most]{tcolorbox}
\tcbuselibrary{breakable, skins}

\theoremstyle{definition}
\newtheorem{definition}{Definition}[chapter]
\newtheorem{proposition}[definition]{Proposition}
\newtheorem{fact}[definition]{Fact}
\crefname{definition}{Definition}{Definitions}
\crefname{proposition}{Proposition}{Propositions}

\tcolorboxenvironment{definition}{
  enhanced, breakable, boxrule=0pt, frame hidden, sharp corners,
  borderline west={2.5pt}{0pt}{keyblue}, colback=keyblue!4,
  left=6pt, right=4pt, top=3pt, bottom=3pt, before skip=8pt, after skip=8pt}
\tcolorboxenvironment{proposition}{
  enhanced, breakable, boxrule=0pt, frame hidden, sharp corners,
  borderline west={2.5pt}{0pt}{intuitionpurple}, colback=intuitionpurple!4,
  left=6pt, right=4pt, top=3pt, bottom=3pt, before skip=8pt, after skip=8pt}
\tcolorboxenvironment{fact}{
  enhanced, breakable, boxrule=0pt, frame hidden, sharp corners,
  borderline west={2.5pt}{0pt}{examplegray}, colback=examplegray!4,
  left=6pt, right=4pt, top=3pt, bottom=3pt, before skip=8pt, after skip=8pt}

% ---------------------------- Note boxes ------------------------------
% Usage:  \begin{keyidea}[optional subtitle] ... \end{keyidea}
\tcbset{
  notebox/.style={enhanced, breakable, sharp corners=south, arc=2pt,
    boxrule=0.6pt, left=6pt, right=6pt, top=8pt, bottom=4pt,
    fonttitle=\bfseries\sffamily\small, before skip=9pt, after skip=9pt,
    attach boxed title to top left={xshift=6pt, yshift=-2.2mm},
    boxed title style={sharp corners, boxrule=0pt}}
}
\NewDocumentCommand{\NewNoteBox}{m m m}{%
  % #1 env name, #2 label, #3 colour
  \NewTColorBox{#1}{O{}}{notebox, colback=#3!5, colframe=#3!85!black,
    boxed title style={colback=#3!85!black},
    title={#2\ifstrempty{##1}{}{: ##1}}}%
}
\NewNoteBox{keyidea}{Key idea}{keyblue}
\NewNoteBox{intuition}{Intuition}{intuitionpurple}
\NewNoteBox{example}{Example}{examplegray}
\NewNoteBox{caution}{Careful}{cautionred}
\NewNoteBox{examtip}{Exam focus}{examorange}
\NewNoteBox{connection}{Connection}{connectgreen}
\NewNoteBox{aside}{Beyond the slides}{asideteal}

% End-of-lecture summary box
\NewTColorBox{summarybox}{O{Lecture summary}}{enhanced, breakable,
  colback=keyblue!3, colframe=keyblue!85!black, boxrule=0.8pt, arc=2pt,
  fonttitle=\bfseries\sffamily, title={#1}, before skip=12pt, after skip=12pt}

% ---------------------------- Self-check ------------------------------
% \showanswerstrue / \showanswersfalse is set in main.tex.
% Usage:
%   \begin{selfcheck}
%     \qa{Question text}{Answer text}
%   \end{selfcheck}
\newif\ifshowanswers
\showanswerstrue
\NewTColorBox{selfcheck}{O{Self-check questions}}{enhanced, breakable,
  colback=examorange!3, colframe=examorange!85!black, boxrule=0.8pt, arc=2pt,
  fonttitle=\bfseries\sffamily, title={#1},
  before upper={\begin{enumerate}[label=\textbf{Q\arabic*.}, leftmargin=2.2em,
                 itemsep=0.6em]},
  after upper={\end{enumerate}}}
\NewDocumentCommand{\qa}{+m +m}{%
  \item #1%
  \ifshowanswers\par\smallskip
    {\color{answercolor}\textbf{Answer.}}\ #2%
  \fi}

% ---------------------------- Misc helpers ----------------------------
% Slide reference in the margin:  \slides{12--15}
% (put it on the line right after \section / \subsection)
\ExplSyntaxOn
\NewDocumentCommand{\slides}{m}{
  \marginnote{\raggedright\footnotesize\sffamily\color{black!50}
    \tl_if_in:nnTF {#1} {-} {Slides}
      { \tl_if_in:nnTF {#1} {,} {Slides} {Slide} } ~ #1}}
\ExplSyntaxOff
% Keep a heading glued to its first paragraph when \slides is used in
% vertical mode (stock marginnote inserts a page-break point there).
\makeatletter
\renewcommand{\mn@vlap}[1]{%
  \setbox\@tempboxa\vbox to \ht\strutbox{#1\vss}%
  \nobreak\box\@tempboxa\nobreak\vskip-\baselineskip\nobreak}
\makeatother
% Defined term (bold italic):  \term{node embedding}
\NewDocumentCommand{\term}{m}{\textbf{\emph{#1}}}
% Lecture heading:  \lecture{<number>}{<title>}{<date>}
\NewDocumentCommand{\lecture}{m m m}{%
  \setcounter{chapter}{\numexpr #1 - 1\relax}%
  \chapter{#2}%
  {\large\itshape\color{black!65} #3\par}\medskip}

\raggedbottom
